\documentclass[11pt]{article}
\usepackage[a4paper,margin=2.6cm]{geometry}
\usepackage{amsmath,amssymb,amsthm}
\usepackage{booktabs}
\usepackage[colorlinks=true,linkcolor=blue,urlcolor=blue]{hyperref}

\newtheorem{claim}{Claim}
\newtheorem{counterexample}{Counterexample}
\newcommand{\Fq}{\mathbb{F}_q}
\newcommand{\Mat}{\operatorname{Mat}}

\title{Disproof of TLMC Conjecture 00000001854:\\
``Squarefree characteristic polynomial counting'' is not an exact closed form}
\author{TLMC submission}
\date{September 13, 2026}

\begin{document}
\maketitle

\begin{abstract}
We disprove the conjecture that
\[
\#\{M \in \Mat_n(\Fq) : \operatorname{charpoly}(M)\ \text{squarefree}\}
\;=\;
q^{n^2}\!\prod_i (1 - q^{-i^2})
\]
is an \emph{exact closed form}. At the minimal dimension $n = 1$ the true count
is exactly $q$ (the characteristic polynomial of a $1 \times 1$ matrix is the
degree-one polynomial $X - a$, which is always squarefree), whereas the
conjectured expression evaluates to $q - 1$ when the product is read over
$i = 1..n$, and to a quantity strictly smaller than $q$ when the product is read
as an infinite product. A second, fully enumerated corroboration at $n = 2$,
$q = 2$ shows a count of $8$ against a formula value of $15/2$, which is not
even an integer.
\end{abstract}

\section{The conjecture}

\begin{claim}[Conjecture 00000001854, literal statement]
\label{conj}
Counting random matrices over $\Fq$:
\[
\#\{M \in \Mat_n(\Fq) : \operatorname{charpoly}(M) \text{ squarefree}\}
\;=\;
q^{n^2}\cdot\prod_i\bigl(1 - q^{-i^2}\bigr)
\]
is an \emph{exact closed form} (squarefree char-poly counting).
\end{claim}

The statement does not specify the index range of the product. We therefore
refute the literal ``exact'' assertion under \emph{both} natural readings:
\begin{itemize}
  \item \textbf{Reading A:} the product runs over $i = 1, \dots, n$ (matching
  the degree of the characteristic polynomial);
  \item \textbf{Reading B:} the product is the infinite product
  $\prod_{i \ge 1}(1 - q^{-i^2})$.
\end{itemize}

\section{Counterexample: \texorpdfstring{$n = 1$}{n = 1}}

\begin{counterexample}[minimal dimension]
\label{ce:n1}
Let $n = 1$. Every element of $\Mat_1(\Fq)$ is a matrix $[a]$ with $a \in \Fq$,
and its characteristic polynomial is
\[
\operatorname{charpoly}([a]) \;=\; X - a .
\]
This is a polynomial of degree $1$. Over any field, a degree-one polynomial $f$
has derivative $f' = 1$, whence
\[
\gcd(f, f') \;=\; \gcd(X - a,\, 1) \;=\; 1,
\]
so $f$ has no repeated irreducible factor, i.e.\ $f$ is squarefree.
Therefore \emph{every} matrix of $\Mat_1(\Fq)$ is counted:
\[
\#\{M \in \Mat_1(\Fq) : \operatorname{charpoly}(M) \text{ squarefree}\} = |\Fq| = q .
\]
\end{counterexample}

Against this, under Reading~A the conjectured formula gives
\[
q^{1}\cdot\Bigl(1 - q^{-1}\Bigr) \;=\; q - 1 \;\neq\; q ,
\]
so Claim~\ref{conj} fails for \emph{every} prime power $q \ge 2$ already at
$n = 1$. Explicit values (all confirmed by the brute-force enumeration of the
accompanying script \texttt{reproduce.py}, which implements polynomial
arithmetic and the squarefreeness test $\gcd(f, f') = 1$ over $\Fq$ for prime
$q$ and over $\mathbb{F}_4 = \mathbb{F}_2[\omega]/(\omega^2 + \omega + 1)$):

\begin{table}[h]
\centering
\begin{tabular}{lccc}
\toprule
$q$ & true count (enumerated) & Reading A: $q(1 - q^{-1})$ & Reading B: $q\prod_{i\ge1}(1-q^{-i^2})$ \\
\midrule
$2$ & $2$ & $1$   & $\approx 0.935655$ \\
$3$ & $3$ & $2$   & $\approx 1.975208$ \\
$4$ & $4$ & $3$   & $\approx 2.988270$ \\
$5$ & $5$ & $4$   & $\approx 3.993598$ \\
\bottomrule
\end{tabular}
\caption{The counterexample $n = 1$. Reading B is computed with exact
rational arithmetic (truncated at $i \le 10$; the omitted tail is bounded by
$q^{-121}$, far below any rounding effect).}
\end{table}

Under Reading~B the failure is structural: every factor $(1 - q^{-i^2})$ lies
strictly between $0$ and $1$, so
$q\prod_{i\ge1}(1 - q^{-i^2}) < q$, while the left-hand side is exactly $q$.
Under either reading, the ``exact closed form'' is false at the smallest
possible dimension.

\section{Second corroboration: \texorpdfstring{$n = 2$, $q = 2$}{n = 2, q = 2}}

We enumerated all $2^{2 \cdot 2} = 16$ matrices of $\Mat_2(\mathbb{F}_2)$,
computed $\operatorname{charpoly}(M) = \det(XI - M)$ with exact polynomial
arithmetic over $\mathbb{F}_2$, and tested squarefreeness via
$\gcd(f, f')$ having degree $0$. The count of matrices with squarefree
characteristic polynomial is
\[
\#\{M \in \Mat_2(\mathbb{F}_2) : \operatorname{charpoly}(M) \text{ squarefree}\} = 8 ,
\]
(the admissible characteristic polynomials being $X^2 + X = X(X+1)$, attained
by $6$ matrices, and the irreducible $X^2 + X + 1$, attained by $2$ matrices;
the remaining $8$ matrices have characteristic polynomial $X^2$ or
$X^2 + 1 = (X+1)^2$, both non-squarefree).

Reading A of the conjectured formula yields instead
\[
2^{4}\cdot\Bigl(1 - 2^{-1}\Bigr)\Bigl(1 - 2^{-4}\Bigr)
= 16 \cdot \frac12 \cdot \frac{15}{16} = \frac{15}{2},
\]
which is \emph{not an integer} --- impossible for an exact count of a finite
set. Reading B yields $\approx 7.485237 \neq 8$. This corroborates
Counterexample~\ref{ce:n1} at the next dimension.

\section{Boundary of the refutation}

Because the index range of $\prod_i(1 - q^{-i^2})$ is unspecified in the
original statement, some charity is required to read it at all. Our refutation
targets the literal \emph{exactness} assertion and is insensitive to the
choice between Readings A and B; both fail, and in both directions (too small
by one at $n = 1$ under Reading A; strictly below the true count under
Reading B; non-integral at $n = 2$ under Reading A). If the product was
intended relative to some further, different normalization (for instance a
probability $\prod_i (1 - q^{1 - i^2})$-type density), then the statement as
written is ambiguous rather than exact; we refute the statement in the form
quoted in Section~\ref{conj}. For reference, the classical exact count of
\emph{monic squarefree polynomials} of degree $n \ge 2$ over $\Fq$ is
$q^n - q^{n-1}$, a statement about polynomials rather than matrices, which
does not coincide with the conjectured expression either.

\section{Machine verification}

The package accompanying this document contains:

\begin{enumerate}
  \item \texttt{reproduce.py}: a dependency-free Python script that
  (i) implements $\Fq$ for prime $q \in \{2,3,5\}$ and
  $\mathbb{F}_4 = \mathbb{F}_2[\omega]/(\omega^2+\omega+1)$, polynomial long
  division, Euclidean $\gcd$, and the squarefreeness test
  $\deg\gcd(f, f') = 0$; (ii) enumerates \emph{all} matrices of
  $\Mat_1(\Fq)$ for $q = 2,3,4,5$ and of $\Mat_2(\mathbb{F}_2)$; and
  (iii) asserts every inequality reported above, terminating successfully
  only if the conjecture fails as described.
  \item \texttt{lean4/}: a Lean~4 project (toolchain
  \texttt{leanprover/lean4:v4.33.1}, no external dependencies) formalizing the
  numeric content of the $n = 1$, $q = 2$ counterexample with
  \emph{zero axioms and zero \texttt{sorry}}: the enumerated count $2$ and the
  formula value $1$ (encoded as an exact fraction over kernel-computable
  natural-number arithmetic, since core \texttt{Rat} does not reduce inside
  the kernel) are computed by \texttt{rfl} and proved distinct by
  \texttt{decide}. Since Lean core does not ship a polynomial library, the
  squarefreeness of $X - a$ is encoded as the (mathematically justified, see
  Counterexample~\ref{ce:n1}) statement that every $a \in \mathbb{F}_2$ passes
  the squarefreeness predicate; the mathematical argument is the one of
  Section~2 and is verified generically by \texttt{reproduce.py} via
  $\gcd(f, f')$.
\end{enumerate}

\section{Conclusion}

Conjecture 00000001854 is \textbf{false}: the expression
$q^{n^2}\prod_i(1 - q^{-i^2})$ is not an exact closed form for the number of
$n \times n$ matrices over $\Fq$ with squarefree characteristic polynomial.
The correct value at $n = 1$ is $q$, not $q - 1$, and already at $n = 2$,
$q = 2$ the formula produces the non-integer $15/2$ against the true count $8$.

\end{document}
